\subsection{正则元}

由 §2 第 2 节定理 2 的推论 4，正则元 \(g\)（\(G\) 的）可由下列性质之一刻画：

\begin{enumerate}
\def\labelenumi{\alph{enumi})}
\item
  g 中被 Ad g 固定的子代数是一个嘉当子代数。
\item
  \(Z(g)_{0}\) 是 G 的一个极大环面。
\end{enumerate}

\(\mathrm{G}\) 的正则元所成的集合在 \(\mathrm{G}\) 中是开且稠密的。

在本段余下部分，我们用 \(G_{r}\)（相应地 \(T_{r}\)）表示 G（相应地 T）中在 G 中正则的点所成的集合。G 的元素 g 属于 \(T_{r}\) 当且仅当 \(Z(g)_{0}\) 等于 T；\(G_{r}\) 的每个元素都共轭于 \(T_{r}\) 的某个元素（§2，第 2 节，定理 2）。

T 的元素 t 属于 \(T_{r}\) 当且仅当 \(t^{\alpha} \neq 1\) 对每个根 \(\alpha \in \mathrm{R}(\mathrm{G}, \mathrm{T})\) 成立；因此，\(T - T_{r}\) 是诸子环面 Ker \(\alpha\)（\(\alpha\) 取遍 \(\mathrm{R}(\mathrm{G}, \mathrm{T})\)）的并。

命题 1. 令 \(n = \dim G\)。存在一个紧实解析流形 \(V\)（维数为 \(n - 3\)），以及一个解析映射 \(\varphi : V \to G\)，其像为 \(G - G_r\)。

设 \(\alpha\in\mathrm{R}(\mathrm{G},\mathrm{T})\)；令 \(\mathrm{V}_{\alpha}=(\mathrm{G}/\mathrm{Z}(\mathrm{Ker}\alpha))\times(\mathrm{Ker}\alpha)\)，并用 \(\varphi_{\alpha}\) 表示从 \(V_{\alpha}\) 到 G 的态射，使得对所有 \(g\in G\) 与所有 \(t\in\operatorname{Ker}\alpha\)，有 \(\varphi_{\alpha}(\bar{g},t) = gtg^{-1}\)（用 \(\bar{g}\) 表示 \(g\) 模 \(\mathrm{Z}(\mathrm{Ker}\alpha)\) 的陪集）。则 \(\mathrm{V}_{\alpha}\) 是一个维数为

\[
\begin{array}{c} \dim V _ {\alpha} = \dim G - \dim Z (\operatorname{Ker} \alpha) + \dim \operatorname{Ker} \alpha \\ = n - (\dim T + 2) + (\dim T - 1) = n - 3 \end{array}
\]

的紧实解析流形（§4，第 5 节，定理 1）；\(\varphi_{\alpha}\) 是实解析流形的态射，且 \(\varphi_{\alpha}\) 的像由 G 中共轭于 \(\mathrm{Ker} \alpha\) 中元素的元素组成。此时只需取 V 为诸流形 \(V_{\alpha}\) 之和，并取 \(\varphi\) 为诱导 \(\varphi_{\alpha}\)（在每个 \(V_{\alpha}\) 上）的态射。

注. 称 G 的元素 g 为强正则的，如果 \(Z(g)\) 是 G 的一个极大环面。若 \(g \in T\)，则 g 强正则当且仅当 \(w(g) \neq g\) 对 \(W_{G}(T)\) 的每个非恒等元素 w 成立（§4，第 7 节，命题 14）。于是，强正则元所成的集合是 G 的一个稠密开子集（§2，第 5 节，命题 5 的推论 2）。

